\begingroup
\renewenvironment{Proof}{\par}{\par\nobreak\hfill$\blacksquare$\par\vspace{0.1cm}}
\begin{Proof}
现在估计主部. 断言每个函数 $a_{\sigma\gamma}$ 都可延拓为整个 $\mathbb{R}^n$ 上的连续函数. (对于 Lipschitz 区域, 已知这一点成立, 因为它们对任意 $k$ 都具有 $k$ 阶延拓性质. 事实上, 由 Tietze 延拓定理 (参阅拓扑学教材), 对任意区域 $\Omega$ 都成立.)

现在设 $\Omega\prime$ 为任意有界开区域, 使得 $\Omega$ 紧包含于 $\Omega\prime$. 由于每个延拓后的 $a_{\sigma\gamma}$ ($|\sigma|=|\gamma|=k$) 在 $\Omega\prime$ 上一致连续, 存在非减的\emph{连续模}函数 $\omega:[0,\infty)\to[0,\infty)$, 满足
\begin{equation}\label{eq:garding-common-continuity-modulus-at-zero}
0=\omega(0)=\lim_{\delta\to0+}\omega(\delta)
\end{equation}
且
\begin{equation}\label{eq:garding-common-coefficient-continuity-modulus}
|a_{\sigma\gamma}(\bm{x})-a_{\sigma\gamma}(\bm{y})|\leq\omega(|\bm{x}-\bm{y}|)
\end{equation}
对每个 $|\sigma|=|\gamma|=k$ 及每个 $\bm{x},\bm{y}\in\Omega\prime$ 成立.

现在对某个 $\bm{x}_0\in\Omega$, 令 $B=B(\bm{x}_0,\delta)$. 稍后再选取 $\delta>0$, 目前仅假设它足够小, 使得 $B\subset\Omega\prime$. 估计 $I_1$ 的第一步是在 $u\in H_0^k(B)$ 的情形下进行估计. 此时有
\begin{equation}\label{eq:garding-frozen-principal-coefficient-splitting}
I_1=I_{11}+I_{12},
\end{equation}
其中
\begin{equation}\label{eq:garding-frozen-principal-part}
I_{11}:=\sum_{|\sigma|=|\gamma|=k}\int_{\mathbb{R}^n}a_{\sigma\gamma}(\bm{x}_0)D^\gamma uD^\sigma u\,\mathrm{d}\bm{x},
\end{equation}

\begin{equation}\label{eq:garding-coefficient-variation-part}
I_{12}:=\sum_{|\sigma|=|\gamma|=k}\int_B[a_{\sigma\gamma}(\bm{x})-a_{\sigma\gamma}(\bm{x}_0)]D^\gamma uD^\sigma u\,\mathrm{d}\bm{x}.
\end{equation}
(注意, 在 $I_{11}$ 的定义中, 假设 $u$ 以 $0$ 延拓到整个 $\mathbb{R}^n$.)

利用 \eqref{eq:garding-common-coefficient-continuity-modulus} 和 Hölder 不等式, 很容易得到 $I_{12}$ 的估计:
\begin{equation}\label{eq:garding-coefficient-variation-modulus-bound}
|I_{12}|\leq\omega(\delta)\|u\|_{k,2}^2.
\end{equation}
为估计 $I_{11}$, 使用 Fourier 变换\[\begin{aligned}
I_{11}&=\sum_{|\sigma|=|\gamma|=k}a_{\sigma\gamma}(\bm{x}_0)\int_{\mathbb{R}^n}D^\gamma u(\bm{x})\overline{D^\sigma u}(\bm{x})\,\mathrm{d}\bm{x}\\
&=\sum_{|\sigma|=|\gamma|=k}a_{\sigma\gamma}(\bm{x}_0)\int_{\mathbb{R}^n}\widehat{D^\gamma u}(\bm{\xi})\overline{\widehat{D^\sigma u}}(\bm{\xi})\,\mathrm{d}\bm{\xi}\\
&=\sum_{|\sigma|=|\gamma|=k}\int_{\mathbb{R}^n}a_{\sigma\gamma}(\bm{x}_0)(i\bm{\xi})^\gamma(-i\bm{\xi})^\sigma|\widehat u|^2\,\mathrm{d}\bm{\xi}\\
&=\sum_{|\sigma|=|\gamma|=k}\int_{\mathbb{R}^n}\bm{\xi}^\sigma a_{\sigma\gamma}(\bm{x}_0)\bm{\xi}^\gamma|\widehat u|^2\,\mathrm{d}\bm{\xi}\\
&\geq\theta\int_{\mathbb{R}^n}|\bm{\xi}|^{2k}|\widehat u|^2\,\mathrm{d}\bm{\xi}.
\end{aligned}\]最后一个不等式使用了一致椭圆性条件. 接着, 利用 Theorem~\ref{thm:sobolev-fourier-weighted-homeomorphism} 得到\[\begin{aligned}
I_{11}&\geq\theta\int_{\mathbb{R}^n}(1+|\bm{\xi}|^{2k})|\widehat u|^2\,\mathrm{d}\bm{\xi}-\theta\int_{\mathbb{R}^n}|\widehat u|^2\,\mathrm{d}\bm{\xi}\\
&\geq\overline C\|u\|_{k,2}^2-\theta\|u\|_2^2
\end{aligned}\]其中 $\overline C>0$ 只依赖于 $\Omega\prime$.

现在结合 $I_{11}$ 和 $I_{12}$ 的估计来估计 $I_1$. 此时假设 $\delta$ 充分小, 使得 $\omega(\delta)\leq\overline C/2$. 那么有
\begin{equation}\label{eq:garding-ball-local-principal-part-bound}
\begin{aligned}
I_1&\geq I_{11}-|I_{12}|\\
&\geq\overline C\|u\|_{k,2}^2-\theta\|u\|_2^2-\omega(\delta)\|u\|_{k,2}^2\\
&\geq\dfrac{\overline C}{2}\|u\|_{k,2}^2-\theta\|u\|_2^2.
\end{aligned}
\end{equation}
现在继续对一般 $u\in H_0^k(\Omega)$ 估计 $I_1$. 基本思想是利用单位分解将 $u$ 分开, 从而能够使用前面的估计.

首先用有限个球覆盖 $\overline\Omega$:
\begin{equation}\label{eq:garding-localization-finite-ball-cover}
B_i:=B(\bm{x}_i,\delta_i),\quad i=1,\ldots,M,
\end{equation}
其中 $\bm{x}_i\in\Omega$, $\delta_i>0$, 如前面的估计那样选取, 使得 $\overline B_i\subset\Omega\prime$. 现在设 $\psi_i$ 为 $\overline\Omega$ 上从属于覆盖 $\overline B_i$ 的单位分解. 然后令
\begin{equation}\label{eq:garding-square-normalized-partition}
\phi_i(\bm{x}):=\left(\dfrac{\psi_i^2(\bm{x})}{\sum_{j=1}^M\psi_j^2(\bm{x})}\right)^{1/2}.
\end{equation}
则有\begin{enumerate}
\item $0\leq\phi_i(\bm{x})\leq1$,
\item $\phi_i\in C^\infty(B_i\cap\overline\Omega)$,
\item $\sum_{i=1}^M\phi_i^2(\bm{x})=1$ 对每个 $\bm{x}\in\Omega$, 且
\item $u_i:=u\phi_i\in H_0^k(B_i)$.
\end{enumerate}利用这些性质可写为\[\begin{aligned}
I_1&=\sum_{|\sigma|=|\gamma|=k}\int_\Omega a_{\sigma\gamma}(\bm{x})D^\sigma uD^\gamma u\,\mathrm{d}\bm{x}\\
&=\sum_{i=1}^M\sum_{|\sigma|=|\gamma|=k}\int_\Omega a_{\sigma\gamma}(\bm{x})\phi_iD^\sigma u\phi_iD^\gamma u\,\mathrm{d}\bm{x}\\
&=\sum_{i=1}^M\sum_{|\sigma|=|\gamma|=k}\int_\Omega a_{\sigma\gamma}(\bm{x})D^\sigma(\phi_i u)D^\gamma(\phi_i u)\,\mathrm{d}\bm{x}\\
&\quad+\sum_{i=1}^M\sum_{|\sigma|=|\gamma|=k}\int_\Omega a_{\sigma\gamma}(\bm{x})[\phi_iD^\sigma u-D^\sigma(\phi_i u)]\phi_iD^\gamma u\,\mathrm{d}\bm{x}\\
&\quad+\sum_{i=1}^M\sum_{|\sigma|=|\gamma|=k}\int_\Omega a_{\sigma\gamma}(\bm{x})[\phi_iD^\gamma u-D^\gamma(\phi_i u)]D^\sigma(\phi_i u)\,\mathrm{d}\bm{x}\\
&\geq\sum_{i=1}^M\sum_{|\sigma|=|\gamma|=k}\int_\Omega a_{\sigma\gamma}(\bm{x})D^\sigma u_iD^\gamma u_i\,\mathrm{d}\bm{x}\\
&\quad-C\|u\|_{k,2}\|u\|_{k-1,2}.
\end{aligned}\]现在对每个 $u_i\in H_0^k(B_i)$ 使用前面的估计, 得到\[\begin{aligned}
I_1&\geq\dfrac{\overline C}{2}\sum_{i=1}^M\|u_i\|_{k,2}^2-C\|u\|_2^2-C\|u\|_{k,2}\|u\|_{k-1,2}\\
&=\dfrac{\overline C}{2}\sum_{i=1}^M\sum_{|\alpha|\leq k}\int_\Omega|D^\alpha(\phi_i u)|^2\,\mathrm{d}\bm{x}\\
&\quad-C[\|u\|_2^2+\|u\|_{k,2}\|u\|_{k-1,2}]\\
&\geq\dfrac{\overline C}{2}\sum_{|\alpha|\leq k}\int_\Omega\sum_{i=1}^M(\phi_i^2)|D^\alpha u|^2\,\mathrm{d}\bm{x}\\
&\quad-C[\|u\|_2^2+\|u\|_{k,2}\|u\|_{k-1,2}]\\
&=\dfrac{\overline C}{2}\|u\|_{k,2}^2-C[\|u\|_2^2+\|u\|_{k,2}\|u\|_{k-1,2}]
\end{aligned}\]因而, 利用 \eqref{eq:garding-elementary-young-product-inequality} 和 \eqref{eq:ehrling-compact-intermediate-norm-bound}, 得到
\begin{equation}\label{eq:garding-global-principal-part-lower-bound}
I_1\geq\dfrac{\overline C}{4}\|u\|_{k,2}^2-C\|u\|_2^2.
\end{equation}
最后将此估计与 \eqref{eq:garding-final-lower-order-estimate} 结合, 得到\[\begin{aligned}
B[u,u]&=I_1+I_2\\
&\geq\left(\dfrac{\overline C}{4}-\varepsilon\right)\|u\|_{k,2}^2-C(\varepsilon)\|u\|_2^2\\
&:=c_3\|u\|_{k,2}^2-\lambda_G\|u\|_2^2,
\end{aligned}\]其中在定义 $c_3$ 和 $\lambda_G$ 时, 已选取充分小的 $\varepsilon$, 例如 $\varepsilon=\overline C/8$.
\end{Proof}
\endgroup
